%% kapa: exam=2 round=33 year=2022 period=2 subject=appraisal_theory subject_ko=감정평가이론 seq=2 number=2 points=30
%% source: Q-Net 감정평가사 기출문제 (cst00309) artlSeq=5210838 file=raw/1008/5210838/2022년 제33회 감정평가사 2차시험 2교시 문제지(감정평가이론).pdf
%% conversion: kordoc 4.24.51 -> tools/gen2.py (auto); review=proofread-gpt-astra; flags=none
\begin{longquestion}{【문제 2】}
소득접근법에서 자본환원율을 결정하는 방법이다. 다음 물음에 답하시오. (30점) \par
\medskip\noindent 물음 1) 투자결합법(band of investment method)의 2가지 유형을 구분하여 쓰고, 엘우드(Ellwood)법을 비교 설명하시오. (20점) \par
\medskip\noindent 물음 2) 자본환원율(capitalization rate)의 조정이 필요한 이유와 조정 방법을 설명하시오. (10점) \par
\end{longquestion}
